\section{Simulation}\label{sec:simulation}

\subsection{Algorithm}
The process is a continuous-time Markov chain with finitely many possible events, so it can be
simulated exactly by drawing the waiting time until the next event from an exponential distribution
with the total rate and then choosing the event with probability proportional to its rate
\citep{gillespie1977,bortz1975}. The number of possible events grows quadratically with the number of
entities, but in sparse networks most of them are spontaneous establishments of links between
entities without any adjacent angle, whose rates are equal within an event type. \tricl\ therefore
schedules three kinds of events:
\begin{itemize}
\item \emph{Individual events}: all terminations of existing links and all establishments with at
  least one adjacent angle of an influencing type. Their rates are the leaves of a complete binary
  \emph{sum tree}, whose inner nodes hold the sums of their children. Updating a rate costs
  $O(\log n)$ for $n$ scheduled events; the root holds the total, and choosing an event with
  probability proportional to its rate is a descent from the root.
\item \emph{Summary events}, one per type of spontaneous establishment, which stand for all pairs of
  entities of the two types. A summary event's leaf holds the upper bound
  $a^0_\tau p^{\max}_\tau N_1 N_3$, with $N_1, N_3$ the numbers of entities of the two types and
  $p^{\max}_\tau$ the largest success probability a covered pair can have. When it is chosen, a pair is
  drawn uniformly; the attempt is discarded if the pair is identical, already linked or scheduled
  individually, and otherwise succeeds with probability $f(B)/p^{\max}_\tau$. This is thinning in the
  sense of \citet{lewis1979}: proposals arrive at the upper-bound rate and are accepted with the ratio
  of the true to the bounding rate, which yields exactly the true rate for every pair.
\item \emph{Immediate events} (infinite rates) are kept in a separate list. While it is not empty,
  one of its events is drawn uniformly and performed without advancing time.
\end{itemize}
After an event on the pair $(e_1, e_3)$, only events whose angle counts change need new rates: those
of pairs $(e_1', e_3)$ with a link between $e_1'$ and $e_1$, and of pairs $(e_1, e_3')$ with a link
between $e_3$ and $e_3'$, for the relationship types that occur in rules. The cost per event is thus
$O\bigl((d_1 + d_3) \log n\bigr)$ in the degrees $d_1, d_3$ of the endpoints. Adjacent angles are found
by intersecting the sorted link sets of the two endpoints, with a galloping search when one set is
much larger than the other, so that the intersection costs time proportional to the smaller set
times the logarithm of the larger one; a hub entity linked to all others (the usual encoding of
states) therefore does not make every event cost time linear in the number of entities. An earlier version scheduled events by
tentative times in an ordered map \citep{gibson2000}; the direct method needs two random numbers per
event instead of one per rate change and gives the exact total rate at all times, which the
likelihood (Sec.~\ref{sec:inference}) needs.

\subsection{Tests}
The test suite runs in continuous integration in optimized, debug and sanitizer builds. It checks
the sigmoid (\TestSigmoidChecks\ unit checks of values, limits, monotonicity and derivatives), the
exact log-likelihood against a closed form for a two-entity model, analytic gradients against finite
differences, that replaying a simulated event sequence reproduces the simulation's log-likelihood,
and, in debug builds, the consistency of every scheduled rate with a recomputation from the
network. Replacing the scheduler left the distribution of all reported statistics unchanged within
sampling error (between 60 and 400 seeds per model; all $|z| < 1.6$).

\subsection{Performance}
Table~\ref{tab:benchmark} reports throughput on one core. Sparse models run at
$10^4$--$10^5$ events per second, and a hundredfold increase in the number of agents of the threshold
model reduces the throughput only about fourfold. In dense models every event changes the rates of
hundreds of others, so the throughput is limited by the number of rate updates. Memory use is a few
tens of megabytes for these models.

\begin{table}[t]
\caption{Throughput of \tricl\ on one core of a cloud virtual machine (best of three runs of
$2\times 10^4$ events, initialization excluded). Helfmann: threshold model of
Fig.~\ref{fig:frames} with mean degree 10; 3-blocks: the triadic-closure model of
Sec.~\ref{sec:model}, mean degree about 100; SIR: epidemic with social distancing on a random
geometric graph.}
\label{tab:benchmark}
\begin{tabular}{llrrr}
\toprule
model & agents & events & initialisation [s] & events per second \\
\midrule
Helfmann & 100 & 20\,000 & 0.07 & 245\,400 \\
Helfmann & 1\,000 & 20\,000 & 0.14 & 164\,600 \\
Helfmann & 10\,000 & 20\,000 & 2.25 & 62\,200 \\
3-blocks (dense) & 300 & 20\,000 & 1.08 & 2\,400 \\
SIR with distancing & 1\,000 & 7\,155 & 0.25 & 33\,000 \\
\bottomrule
\end{tabular}

\end{table}
